\documentclass[10pt,letterpaper]{article}
\usepackage{cornell-notes}

\title{Clause 22.10.07.07: Class Template Negate}
\author{}
\date{}
\setDocTitle{Clause 22.10.07.07: Class Template Negate}
\setDocSubtitle{C++ 2024}
\setDocOwner{C++ 2024 Cornell Notes}
\setDocDate{\today}
\setCornellCollection{C++ 2024 Cornell Notes}
\setCornellUnitType{Clause}
\setCornellUnitNumber{22.10.07.07}
\setCornellUnitTitle{Class Template Negate}
\hypersetup{pdftitle={Clause 22.10.07.07: Class Template Negate},pdfsubject={Cornell notes},pdfauthor={}}

\begin{document}
\maketitle
\begin{CornellOverviewBox}{Overview}
\textbf{Classification:} Standard library - Normative\\[2pt]
\textbf{Learning objective:} Understand the rules, terminology, and semantic consequences associated with Class template negate.
\end{CornellOverviewBox}\begin{CornellSummaryBox}{Summary}
{\sffamily\bfseries\color{Accent} Summary}\par\vspace{3pt}Section 22.10.7.7 focuses on Class template negate. Apply the specified interface together with its mandates, effects, result conditions, exception behavior, and complexity constraints. When applying it, preserve the distinction between mandatory requirements, recommendations, permissions, and behavior deliberately left undefined, unspecified, or implementation-defined.
\end{CornellSummaryBox}

\end{document}
